\subsection{Boussinesq\_map.txt}
\index{Boussinesq coefficient}
%
\subsubsection{Purpose}
The Boussinesq momentum distribution coefficient is a cell attribute. Two options are available for specifying this coefficient:
\begin{enumerate}
    \item $\beta$ is uniform over the entire domain, in which case providing a single value is sufficient,
    \item $\beta$ is spatially variable, in which case its value must be specified on a cell-by-cell basis.
\end{enumerate}
%
\subsubsection{Format description}
%
File format for a uniform Boussinesq coefficient. \texttt{\textbf{Bold}} : keyword, must appear as such in the file. \texttt{<Tab>}: tabulation character. \\
    \begin{tabular}{p{5cm} | p{10 cm}}
    \hline
    Field(s) & Comment(s) \\
    \hline  \hline
    \texttt{\textbf{Unif} <Tab> \textbf{1}} & The value 1 for the flag indicates that $\beta$ is uniform \\ \hline
    \texttt{\textbf{==== Default Param}} & Section separator line, leave unchanged \\  \hline
    \texttt{\textbf{Boussinesq}} & Comment line, leave unchanged \\  \hline
    \texttt{beta}     & Uniform value \\ 
    \texttt{\textbf{==== Distrib}} & Section separator line, leave unchanged \\  
    \hline
    \end{tabular}
%
\vspace{10mm} \\
%
File format for a spatially variable Boussinesq coefficient. \texttt{\textbf{Bold}} : keyword, must appear as such in the file. \texttt{<Tab>}: tabulation character.  \\
    \begin{tabular}{p{5cm} | p{10 cm}}
    \hline
    Field(s) & Comment(s) \\
    \hline  \hline
    \texttt{\textbf{Unif} <Tab> \textbf{0}} & The zero value indicates that $\beta$ is spatially variable \\ \hline
    \texttt{\textbf{==== Default Param}} & Section separator line, leave unchanged \\  \hline
    \texttt{\textbf{Boussinesq}} & Comment line, leave unchanged \\  \hline
    \texttt{beta}     & Uniform value, will be ignored because the flag is 0 \\ 
    \hline
    \texttt{\textbf{==== Distrib}} & Section separator line, leave unchanged \\  \hline
    \texttt{beta[i]} & $\beta$ value for each cell. Running index $i$ is the index of the cell in the mesh. One value per line.  \\ 
    \hline
    \end{tabular}
%%%
\subsubsection{Sample files}
\paragraph*{Example\#1.} $\beta$ is uniform over the domain.
\begin{verbatim}
Unif	1
==== Default Param
Boussinesq
2.
\end{verbatim}
%
\paragraph*{Example\#2.} $\beta$ is variable over the domain, the mesh is made of 3 cells.
\begin{verbatim}
Unif	0
==== Default Param
Boussinesq
2.
==== Distrib
1.2
1.1
3.5
\end{verbatim}